\section{A checked uniform scheduler and seed proof}
\label{sec:seed-rank}

The current production Rocq compiler proves three related progress facts
separately: it drains a positive proof-list pointer, reaches successive
candidates from the zero start, and repeats the candidate argument from the
bootstrap seed.  We replace these inductions by one numerical rank, retaining
the exact command semantics and the existing public endpoints.

For a chosen candidate number $N$ and state $s$ with $\mathrm{np}(s)\leq N$, set
\[
 R_N(s)=(N-\mathrm{np}(s))(N+1)+\mathrm{pl}(s).
\]
The existing \texttt{step\_pl} theorem describes every successful enumerator
step.  If $\mathrm{pl}(s)>0$, the candidate counter is unchanged and the
proof-list pointer decreases strictly.  If $\mathrm{pl}(s)=0$, then
$\mathrm{np}(s')=\mathrm{np}(s)+1$ and
$\mathrm{pl}(s')\leq\mathrm{np}(s)$.  Unless the desired boundary
$\mathrm{pl}(s)=0$, $\mathrm{np}(s)=N$ has already been reached, the latter
case has $\mathrm{np}(s)<N$.  Consequently
\[
 R_N(s')-R_N(s)
 \leq \mathrm{np}(s)-(N+1)<0.
\]
The next counter remains at most $N$ in both cases.

\paragraph{Checked bounded-search theorem.}
The new \texttt{find\_candidate\_bounded} proves that there exist a natural
number $n\leq R_N(s)$ and a state $t$ such that the existing run function
executes $n$ successful steps from $s$ to $t$, and either $t$ takes a halting
step or $t$ is the boundary for candidate $N$.  This is proved by well-founded
induction on $R_N(s)$.  The halting transition itself, when counted, takes one
additional step; the theorem uses the repository's convention of witnessing
the last nonhalting configuration.

\paragraph{Completeness from every initial state.}
The existing derivation compiler supplies a command word $L$ with contract
\texttt{blk 1 L} whenever $\mathrm{thm222}(1)$ holds.  The already proved
\texttt{program\_halts} theorem states that
\texttt{START Y :: L ++ SUFFIX} halts from every working-register tuple.
Taking $Y=\mathrm{np}(s)$, its exact candidate number is
\[
 N=P\bigl(3,P(Y,P(1,\mathrm{enc\_tail}(L\mathbin{+\!+}\mathrm{SUFFIX})))\bigr)
 \geq Y.
\]
The inequality follows directly from the retained Cantor-pairing inequalities.
The rank theorem reaches this candidate or an earlier halt.  The resulting
new theorem is
\[
 \forall s,\qquad
 \mathrm{thm222}(1)\ \Longrightarrow\ \mathrm{halts\_from}(s).
\]
It assumes neither a zero proof-list pointer nor a particular candidate
counter or working-register value.

For the converse we retain precisely the existing invariant
\[
 \mathrm{Inv}(s)=
 \mathrm{thm222}(\mathrm{tw}(s))\ \land\
 \mathrm{thm222}(\mathrm{h}(s))\ \land\
 \mathrm{goodstack}(\mathrm{ws}(s)).
\]
The checked \texttt{enum\_from\_iff} theorem states
\[
 \forall s,\qquad \mathrm{Inv}(s)\Longrightarrow
 \bigl(\mathrm{halts\_from}(s)\iff\mathrm{thm222}(1)\bigr).
\]
Thus every state $\mathrm{mkState}(p,n,0,0,d,0)$ is a valid start, for arbitrary
$p,n,d$.  This gives future bootstrap proofs a more flexible interface.  It
does not prove a shorter physical bootstrap trace.  The invariant cannot be
removed from the converse: an unjustified falsity code placed directly in the
initial cache can cause an immediate NIP halt.

\paragraph{A separate opacity improvement.}
Adding
\texttt{Local Opaque P wim wal wex weq wel} to
\texttt{L1\_SoundPerm.v} leaves every theorem statement and existing proof
script unchanged.  Those proofs use constructor and inverse lemmas; they do
not require numerical expansion of formula codes.  The observed module
compilation time fell from $43.290$ seconds in the initial run to $0.810$
seconds in a focused trial and $0.665$ seconds in the final fresh build.
These are host-local Coq~8.18 measurements, not a controlled full-project
benchmark or an asserted Rocq~9.1 speedup.

\paragraph{Size and verification.}
The compiler and seed files together shrink from $18\,100$ to $17\,209$
bytes; the seed file shrinks from $59$ to $24$ lines.  Including the opacity
annotation, all three edited files shrink from $22\,105$ to $21\,335$ bytes.
The universal theorem is stronger even though its compiler file alone grows.
All ten required project modules and a typed interface client compile afresh
under Coq~8.18.0, with six closed assumption reports.  Summed compiler time is
$16.996$ seconds; the separate recursive \texttt{coqchk} pass succeeds in
$66.483$ seconds.  Exact sources, patches, compiler commands, hashes and raw
logs accompany the report.

The public types of \texttt{enum\_complete},
\texttt{seedPerm\_complete}, \texttt{seedPerm\_sound}, and
\texttt{seedPerm\_iff\_thm} are preserved and independently elaborated by the
client.  The production \texttt{MainPerm} composition consumes the unchanged
last theorem.  Its full binary-machine-to-ZFC dependency closure was not
rebuilt in this focused check.  Integration should update both the maintained
\texttt{rocqperm} templates and their generated \texttt{K215F} copies, then
rebuild the ordinary production endpoint.
